\clearpage
\fancyhead[L]{\small GAUSS AN BODE.}
\hypertarget{native-419-part-lower}{}\label{native-419-part-lower}
\begin{center}
\Large GAUSS AN BODE.
\end{center}

\begin{center}
\rule{6cm}{0.4pt}

Astronomisches Jahrbuch für 1821. Seite 212..216. Berlin 1818.
\end{center}

\begin{flushright}Göttingen, den 7. Sept. 1818.\end{flushright}

\par\indent Ihrem Wunsche gemäss übersende ich Ihnen, hochverehrter Freund, beigehend abermals eine Ephemeride für die Pallas vom August 1819 bis Juni 1820. Sie ist vom Hrn. \textsc{Dirksen} berechnet, welcher sich hier mit Eifer den mathematischen Wissenschaften widmet. Eben derselbe hat auch die Berechnung der Opposition des laufenden Jahres noch einmal wiederholt nach denselben Elementen, nach welchen Hr. Dr. \textsc{Tittel} gerechnet hatte, und ein etwas abweichendes Resultat gefunden, nemlich:

\begin{center}
\setlength{\tabcolsep}{3pt}
\begin{tabular}{l|c|c}
 & Hr. Dirksen & Hr. Dr. Tittel \\
\hline
Zeit in Göttingen . & \shortstack{Sept. 8\\$15^{\mathrm u}\;17^{\mathrm m}\;5^{\mathrm s}$} & \shortstack{Sept. 8\\$18^{\mathrm u}\;44^{\mathrm m}\;12^{\mathrm s}$} \\
Wahre Länge . . & $345^\circ\;47'\;18''$ & $345^\circ\;55'\;34''$ \\
Geocentrische Breite & $6\;47\;25$ & $6\;47\;11$ \\
\end{tabular}
\end{center}

\par\indent Welches die Entstehung dieses Rechnungsfehlers sei, kann ich nicht sagen; er ist aber auf die Berechnung der Opposition eingeschränkt, da mehrere von Hrn. \textsc{Dirksen} nachgerechnete Oerter der \SourceLeaf{420} Ephemeride vollkommen mit Hrn. Dr. \textsc{Tittels} Angaben übereinstimmen (deren Richtigkeit auch durch meine nachher anzuführenden Beobachtungen bereits bestätigt wird).

\par\indent Seit Ende Aprils d. J. beobachte ich nunmehro an dem trefflichen Meridiankreise von \textsc{Repsold}, welches Instrument Sie ohne Zweifel früher selbst gesehen haben, woran aber jetzt so viele wichtige Verbesserungen angebracht sind, dass es für ein ganz neues gelten kann. Es hat eine ganz neue Theilung erhalten, drei Ablesungen statt der frühern einzigen, neue Zapfenlager von Bergkrystall, die Zapfen selbst sind neu abgedreht und centrirt u.s.w. Verschiedene andere neue Vorrichtungen dazu erwarte ich noch von Hrn. \textsc{Repsold}. Als Probe von der Vortrefflichkeit des Instruments, insofern es Mittagsfernrohr ist, mögen folgende Bestimmungen der Ger. Aufsteig. des Nordsterns dienen; die letzte Columne unter der Ueberschrift Gewicht drücken das Quadrat der Genauigkeit nach Massgabe der Zahl der concurrirenden Fäden aus, wenn die Genauigkeit einer Bestimmung aus einem Fadenantritt in der obern und einem Fadenantritt in der untern Culmination zur Einheit genommen wird. Die Formel dafür ist
\[
\frac{2ab}{a+b}
\]
\noindent wenn, bei einer auf zwei Culminationen gegründeten Bestimmung $a$ und $b$ die Anzahl der beobachteten Antritte in jeder ausdrücken, und
\[
\frac{8abc}{ab+4ac+bc}
\]
\noindent bei einer Bestimmung aus drei auf einander folgenden Culminationen, wo $c$ Fädenantritte in der dritten beobachtet sind. Im ersten Falle setzt man voraus, dass die Pfeiler 12 Stunden unverrückt gestanden; im letztern nur, dass ihre etwanigen Veränderungen während 24 Stunden gleichförmig gewesen sind. Die Zeiten jeder Rectascensionsbestimmung sind von der untern Culmination jedes Tages an gezählt.

\begin{center}
\begin{tabular}{l|r|c|c}
1818 & \shortstack{Beob. Ger. Aufst.\\des Nordsterns} & \shortstack{Correction von\\\textsc{Bessels} Tafel} & Gewicht \\
\hline
Juni 3.25 & $0^{\mathrm u}\,55^{\mathrm m}\,58^{\mathrm s}65$ & $+3^{\mathrm s}93$ & 2.67 \\
4.25 & 58.19 & $+2.80$ & 5.00 \\
5.25 & 58.56 & $+2.45$ & 5.00 \\
6.25 & 58.91 & $+2.22$ & 4.00 \\
8.50 & $0\,\;56\;\;1.14$ & $+2.97$ & 6.40 \\
10.25 & 2.58 & $+3.20$ & 3.75 \\
11.25 & 2.68 & $+2.59$ & 1.67 \\
12.50 & 3.31 & $+2.36$ & 4.29 \\
16.75 & 6.11 & $+2.13$ & 4.44 \\
18.25 & 7.31 & $+2.26$ & 5.00 \\
19.25 & 8.29 & $+2.51$ & 2.86 \\
21.25 & 10.49 & $+3.26$ & 5.00 \\
22.75 & 12.02 & $+3.60$ & 1.33 \\
26.25 & 13.49 & $+2.60$ & 1.60 \\
27.25 & 14.77 & $+3.14$ & 4.44 \\
29.50 & 16.71 & $+3.42$ & 3.43 \\
Juli 7.50 & 22.51 & $+3.31$ & 6.00 \\
8.25 & 22.97 & $+3.22$ & $1\;67$ \\
9.25 & 23.89 & $+3.40$ & 3.75 \\
Sept. 3.75 & 57.44 & $+2.35$ & 2.86 \\
\end{tabular}
\end{center}

\par\indent In sofern man die Correction von \textsc{Bessels} Tafel während dieser Zeit als constant annimmt, wäre dieselbe im Mittel
\[
= +2^{\mathrm s}86
\]

\Needspace{20\baselineskip}\SourceLeaf{421}
\par\indent Noch füge ich meine Beobachtungen des Uranus, des Saturn und der Pallas bei; die Oppositionen der beiden letzten Planeten werden erst in den nächsten Tagen eintreten.

\begin{center}
\emph{Beobachtungen des Uranus.}

\setlength{\tabcolsep}{2.2pt}\begin{tabular}{l|ccc|ccc|ccc}
\hline
1818 & \multicolumn{3}{c|}{Mittlere Zeit} & \multicolumn{3}{c|}{Gerade Aufsteig.} & \multicolumn{3}{c}{Abweichung} \\
\hline
Juni 3 & $12^{\mathrm u}$ & $22^{\mathrm m}$ & $17^{\mathrm s}5$ & $257^\circ$ & $23'$ & $28''3$ & $23^\circ$ & $1'$ & $41''6$ Südl. \\ 
4 & 12 & 18 & 10.9 & & 20 & 47.8 & & 1 & 29.0 \\
5 & 12 & 14 & $4\;6$ & & 18 & 9.9 & & 1 & 14.4 \\
7 & 12 & 5 & 51.6 & & 12 & 54.6 & & 0 & 50.5 \\ 
8 & 12 & 1 & 45.1 & & 10 & 15.0 & & 0 & 38.4 \\
9 & 11 & 57 & 38.6 & & 7 & 36.4 & & 0 & 27.4 \\
10 & 11 & 53 & $32\;1$ & & 4 & 56.5 & & 0 & 12.6 \\
11 & 11 & 49 & 25.7 & & 2 & 18.0 & & 0 & 2.9 \\
12 & 11 & 45 & 19.2 & 256 & 59 & 33.9 & 22 & 59 & 50.9 \\
13 & 11 & 41 & 12.7 & & 57 & 0.4 & & 59 & 40.5 \\
\hline
\end{tabular}
\end{center}

\begin{center}
\emph{Beobachtungen des Saturn.}

\setlength{\tabcolsep}{2.2pt}\begin{tabular}{l|ccc|ccc|ccc}
\hline
Aug. 31 & $12^{\mathrm u}$ & $31^{\mathrm m}$ & $25^{\mathrm s}0$ & $347^\circ$ & $24'$ & $5''2$ & $7^\circ$ & $48'$ & $6''2$ Südl. \\
Sept. 1 & 12 & 27 & 12.0 & & 19 & 49.3 & & 49 & 55.8 \\
3 & 12 & 18 & 47.0 & & 11 & 30.1 & & 53 & 37.4 \\
4 & 12 & 14 & 34.1 & & 7 & 14.1 & & 55 & 27.4 \\
6 & 12 & 6 & 8.7 & 346 & 58 & 50.2 & & 59 & 7.0 \\
\hline
\end{tabular}
\end{center}

\begin{center}
\emph{Beobachtungen der Pallas.}

\setlength{\tabcolsep}{2.2pt}\begin{tabular}{l|ccc|ccc|ccc}
\hline
Aug. 31 & $12^{\mathrm u}$ & $25^{\mathrm m}$ & $8^{\mathrm s}7$ & $345^\circ$ & $49'$ & $45''1$ & $2^\circ$ & $23'$ & $2''4$ Nordl. \\
Sept. 4 & 12 & 6 & 26.8 & & 5 & 5.8 & 1 & 32 & 43.4 \\
\hline
\end{tabular}
\end{center}

\par\indent Die Beobachtungen so lichtschwacher Gestirne wie die Pallas werden sehr dadurch erschwert, dass die äusserst zarten Fäden (die natürlichen Spinnfäden sind mehreremale gespalten), welche übrigens besonders für Tag-Beobachtungen so vortreffliche Dienste leisten, kaum bei derjenigen Beleuchtung sichtbar gemacht werden können, welche jene Gestirne vertragen. Ich werde in Zukunft für dergleichen Zwecke noch besondere dickere Fäden einziehen.

\Needspace{22\baselineskip}\par\indent Die Boobachtungen des Uranus hat Hr. \textsc{Dirksen} mit den \textsc{de Lambre}schen Uranus-Tafeln verglichen und folgende Unterschiede gefunden: 

\begin{center}
\begin{tabular}{l|c|c}
 & Ger. Aufst. & Abweich. \\
\hline
Juni 3 & $-54''1$ & $+10''3$ \\
4 & $-52.1$ & $+10.6$ \\
5 & $-52.9$ & $+12.9$ \\
7 & $-55.8$ & $+12.3$ \\
8 & $-55.9$ & $+12.0$ \\
\end{tabular}
\par\medskip
\begin{tabular}{l|c|c}
 & Ger. Aufst. & Abweich. \\
\hline
Juni 9 & $-56''8$ & $+10''5$ \\
10 & $-56.4$ & $+12.9$ \\
11 & $-57.2$ & $+10.1$ \\
12 & $-52.4$ & $+9.4$ \\
13 & $-58.1$ & $+7.4$ \\
Mittel & $-55.2$ & $+10.8$ \\
\end{tabular}
\end{center}

\Needspace{16\baselineskip}\par\indent Hieraus berechnete Hr. \textsc{Dirksen} die Opposition folgendermassen:

\begin{center}
\begin{tabular}{ll}
\multicolumn{2}{l}{1818 Juni 9. $5^{\mathrm u}\,30^{\mathrm m}\,43^{\mathrm s}$} \\
\multicolumn{2}{l}{Mittlere Zeit in Göttingen.} \\
Wahre Länge . . . . . . . . . . . . . & $258^\circ\;10'\;34''5$ \\
\end{tabular}\par\smallskip
\begin{tabular}{ll}
Geocentrische Breite . . . . . . . . . . . & $0\;4\;12.1$ \\
Heliocentrische . . . . . . . . . . . . . & $0\;3\;58.7$ \\
\end{tabular} $\left.\rule{0pt}{2.2em}\right\}$ Südl.
\end{center}

\noindent und den Fehler jener Tafeln heliocentrisch

\begin{center}
\begin{tabular}{ll}
Länge . . . . . . . . . . . . . . . . & $-45''3$ \\
Breite . . . . . . . . . . . . . . . . & $+14.6$ \\
\end{tabular}
\end{center}

\par\indent Von meinen zahlreichen Fixsternbeobachtungen ist erst ein Theil reducirt, und ich behalte mir die Mittheilung der Resultate auf die Zukunft vor.
